\section{Historical Evolution and Architectural Taxonomy}\label{sec:history}

The mathematical principles formalized in Section~\ref{sec:math} did not emerge in isolation; they evolved across nearly eight decades of cross-pollination among optical physics, cognitive neuroscience, connectionist artificial intelligence, and modern deep sequence modeling.
Figure~\ref{fig:genealogy} maps this intellectual lineage across three historical waves---from 1990s complex-valued Holographic Neural Networks, through algebraic Vector Symbolic Architectures, to contemporary spectral layers and State-Space Models---revealing a deep structural convergence beneath seemingly disparate model families.

\begin{figure}[tbp]
    \centering
    % ==============================================================================
% figures/fig_genealogy.tex
% Historical & Architectural Genealogy (TikZ 4-tier chronological DAG)
% Traces evolution from Optical Holography & 1990s HNNs -> HRR/HDC -> Modern
% CVNNs, Spectral Mixing, SSMs, & Modern Hopfield -> Emerging Hybrid Frontier.
% Requires TikZ libraries: positioning, arrows.meta, shapes.geometric, fit,
% calc, backgrounds; and Okabe-Ito colors + math macros from main.tex.
% ==============================================================================
\begin{tikzpicture}[
    >=Stealth,
    node distance=0.85cm and 0.48cm,
    every node/.style={font=\footnotesize},
    gennode/.style={
        rectangle,
        rounded corners=3pt,
        line width=0.8pt,
        align=center,
        inner sep=4.5pt,
        minimum height=1.18cm
    },
    wave1/.style={
        gennode,
        draw=okabeNavy,
        fill=okabeNavy!7,
        text width=3.95cm
    },
    wave2/.style={
        gennode,
        draw=okabeTeal,
        fill=okabeTeal!8,
        text width=5.95cm
    },
    wave3/.style={
        gennode,
        draw=okabeAmber!85!black,
        fill=okabeAmber!10,
        text width=3.02cm
    },
    wave3alt/.style={
        gennode,
        draw=okabeCoral,
        fill=okabeCoral!8,
        text width=3.02cm
    },
    frontier/.style={
        gennode,
        draw=okabePurple!85!black,
        fill=okabePurple!10,
        text width=6.15cm
    },
    edgeflow/.style={
        ->,
        line width=0.8pt,
        draw=okabeSlate!85
    },
    crossflow/.style={
        ->,
        line width=0.75pt,
        dashed,
        draw=okabeSlate!70
    },
    eraband/.style={
        rounded corners=5pt,
        dashed,
        line width=0.65pt,
        inner sep=6pt
    }
]

% ==============================================================================
% TIER 1: FOUNDATIONS & WAVE 1 (1948 -- 1990s)
% ==============================================================================
\node[wave1] (sutherland) {%
    \textbf{\color{okabeNavy}Sutherland HNN (1990s)}\\[1.5pt]
    {\scriptsize Phasors $\lambda e^{i\theta} \in \C^D$; one-shot\\
    correlation $\mat{X}=\sum \vect{v}_n\vect{k}_n\herm$}%
};

\node[wave1, left=0.55cm of sutherland] (optics) {%
    \textbf{\color{okabeNavy}Optical Holography (1948--68)}\\[1.5pt]
    {\scriptsize Gabor wave interference;\\
    Longuet-Higgins holophone}%
};

\node[wave1, right=0.55cm of sutherland] (hopfield) {%
    \textbf{\color{okabeNavy}Hopfield (1982) \& TPR (1990)}\\[1.5pt]
    {\scriptsize Energy attractors ($\bigo{D^2}$);\\
    Smolensky tensor binding $\vect{r}\otimes\vect{f}$}%
};

% ==============================================================================
% TIER 2: WAVE 2 -- ALGEBRAIC VECTOR SYMBOLIC ARCHITECTURES (1995 -- 2010s)
% ==============================================================================
\coordinate (t1_left_mid) at ($(optics.south)!0.48!(sutherland.south)$);
\node[wave2, below=0.95cm of t1_left_mid] (hrr) {%
    \textbf{\color{okabeTeal}Plate's HRR \& FHRR (1995--2003)}\\[1.5pt]
    {\scriptsize Fixed-$D$ binding ($\cconv \leftrightarrow \hadamard$ on $\mathbb{T}^D$), correlation $\ccorr$,\\
    clean-up memory, recursive symbolic structures}%
};

\node[wave2, right=0.65cm of hrr] (hdc) {%
    \textbf{\color{okabeTeal}Kanerva's HDC \& VSA Taxonomy (2009--)}\\[1.5pt]
    {\scriptsize High-$D$ ($10^4$) binary/bipolar hypervectors (MAP/BSC),\\
    permutation sequences, fault-tolerant bitwise ops}%
};

% ==============================================================================
% TIER 3: WAVE 3 -- MODERN DEEP LEARNING CONVERGENCE (2016 -- PRESENT)
% ==============================================================================
\node[wave3, below=1.05cm of hrr, xshift=-1.75cm] (cvnn) {%
    \textbf{\color{okabeAmber!75!black}CVNNs \& uRNNs}\\[1.5pt]
    {\scriptsize Wirtinger $\nabla$, $\C$-act.,\\
    unitary $\mat{U}\herm\mat{U}=\mat{I}$}%
};

\node[wave3, right=0.45cm of cvnn] (spectral) {%
    \textbf{\color{okabeAmber!75!black}Spectral Mixing}\\[1.5pt]
    {\scriptsize FNet 2D DFT; FFT\\
    long convolutions}%
};

\node[wave3, right=0.45cm of spectral] (ssm) {%
    \textbf{\color{okabeAmber!75!black}Linear Attn.\ \& SSMs}\\[1.5pt]
    {\scriptsize S4/Mamba gated state\\
    $\mat{S}_t{=}\mat{A}_t\mat{S}_{t-1}{+}\vect{v}_t\vect{k}_t\herm$}%
};

\node[wave3alt, right=0.45cm of ssm] (modhop) {%
    \textbf{\color{okabeCoral}Modern Hopfield}\\[1.5pt]
    {\scriptsize Exp.\ energy $\equiv$ softmax\\
    attention ($\bigo{ND}$)}%
};

% ==============================================================================
% TIER 4: EMERGING FRONTIER (SECTIONS 5 -- 6)
% ==============================================================================
\coordinate (t3_left_mid) at ($(cvnn.south)!0.5!(spectral.south)$);
\node[frontier, below=1.0cm of t3_left_mid] (photonics) {%
    \textbf{\color{okabePurple!80!black}Photonic \& Edge Hardware Co-Design (Sec.~\ref{sec:applications})}\\[1.5pt]
    {\scriptsize MZI unitary meshes, diffractive $\text{D}^2\text{NN}$, in-memory RRAM/ASIC HDC}%
};

\coordinate (t3_right_mid) at ($(ssm.south)!0.5!(modhop.south)$);
\node[frontier, below=1.0cm of t3_right_mid] (hybrid) {%
    \textbf{\color{okabePurple!80!black}Hybrid Holographic Foundation Models (Sec.~\ref{sec:potentials})}\\[1.5pt]
    {\scriptsize $\bigo{1}$ holographic KV-caches, neuro-symbolic binding, continual superposition}%
};

% ==============================================================================
% DIRECTED LINEAGE EDGES
% ==============================================================================
% Tier 1 -> Tier 1 / Tier 2
\draw[edgeflow] (optics.east) -- (sutherland.west);
\draw[edgeflow] (optics.south) -- ++(0,-0.25) -| ([xshift=-1.4cm]hrr.north);
\draw[edgeflow] (sutherland.south) -- ++(0,-0.25) -| ([xshift=0.9cm]hrr.north);
\draw[edgeflow] (hopfield.south) -- ++(0,-0.25) -| ([xshift=1.2cm]hrr.north);
\draw[edgeflow] (hopfield.south) -- ++(0,-0.25) -| ([xshift=-0.8cm]hdc.north);

% Horizontal VSA synergy
\draw[<->, okabeTeal, line width=0.8pt] (hrr.east) -- (hdc.west);

% Tier 1/2 -> Tier 3
\draw[edgeflow] (sutherland.south) -- ++(0,-0.25) -| ([xshift=-0.2cm]cvnn.north);
\draw[edgeflow] (hrr.south) -- ++(0,-0.32) -| (spectral.north);
\draw[edgeflow] (hrr.south) -- ++(0,-0.32) -| ([xshift=-0.45cm]ssm.north);
\draw[edgeflow] (hdc.south) -- ++(0,-0.32) -| ([xshift=0.45cm]ssm.north);
\draw[crossflow] (hopfield.east) -- ++(0.35,0) |- (modhop.east);

% Tier 3 -> Tier 4
\draw[edgeflow] (cvnn.south) -- ++(0,-0.32) -| ([xshift=-1.5cm]photonics.north);
\draw[edgeflow] (spectral.south) -- ++(0,-0.32) -| ([xshift=0.8cm]photonics.north);
\draw[edgeflow] (spectral.south) -- ++(0,-0.32) -| ([xshift=-1.8cm]hybrid.north);
\draw[edgeflow] (ssm.south) -- ++(0,-0.32) -| ([xshift=-0.4cm]hybrid.north);
\draw[edgeflow] (modhop.south) -- ++(0,-0.32) -| ([xshift=1.4cm]hybrid.north);

% Direct optical & edge heritage into Tier 4 (routed cleanly outside left margin)
\draw[crossflow] (optics.west) -- ++(-0.35,0) |- (photonics.west);
\draw[crossflow] (hdc.south) -- ++(0,-0.32) -| ([xshift=2.2cm]photonics.north);

% ==============================================================================
% BACKGROUND ERA BANDS
% ==============================================================================
\begin{scope}[on background layer]
    \node[eraband, draw=okabeNavy!40, fill=okabeLightBg,
          fit=(optics)(sutherland)(hopfield),
          label={[anchor=south west, font=\scriptsize\bfseries, text=okabeNavy]north west:Wave 1: Optical Holography, Classical Associative Memories, \& 1990s HNNs}] {};
    \node[eraband, draw=okabeTeal!45, fill=okabeLightBg,
          fit=(hrr)(hdc),
          label={[anchor=south west, font=\scriptsize\bfseries, text=okabeTeal!85!black]north west:Wave 2: Vector Symbolic Architectures (HRR / FHRR \& Hyperdimensional Computing)}] {};
    \node[eraband, draw=okabeAmber!55!black, fill=okabeLightBg,
          fit=(cvnn)(spectral)(ssm)(modhop),
          label={[anchor=south west, font=\scriptsize\bfseries, text=okabeAmber!75!black]north west:Wave 3: Complex Deep Learning, Spectral Mixing, State-Space Recurrences, \& Softmax Attention}] {};
    \node[eraband, draw=okabePurple!55!black, fill=okabeLightBg,
          fit=(photonics)(hybrid),
          label={[anchor=south west, font=\scriptsize\bfseries, text=okabePurple!80!black]north west:Emerging Convergence: Algorithm--Hardware Co-Design \& Hybrid Foundation Models}] {};
\end{scope}

\end{tikzpicture}

    \caption{Historical and architectural genealogy of holographic neural computation. Solid arrows denote direct mathematical or algorithmic inheritance; dashed arrows highlight cross-era physical and architectural bridges linking optical holography to photonic accelerators (Section~\ref{sec:applications}) and classical Hopfield attractors to modern softmax attention (Section~\ref{sec:comparison}).}
    \label{fig:genealogy}
\end{figure}

% ------------------------------------------------------------------------------
\subsection{First Wave: Optical Analogies and 1990s Holographic Neural Networks}\label{sec:history:hnn}
% ------------------------------------------------------------------------------

\subsubsection{Optical and Associative Memory Precursors}

The conceptual genesis of holographic computation traces to Dennis Gabor's invention of optical holography~\cite{GABOR_1948_HOLOGRAPHY}, which demonstrated that a coherent wave field can record both the amplitude and relative phase of an object across an entire physical interference plate.
Because information is distributed globally across the interference pattern, illuminating even a fractured sub-region of the hologram with the original reference beam reconstructs the entire stored image with gracefully degraded resolution rather than localized erasure.
Two decades later, Longuet-Higgins translated this physical mechanism into a mathematical model of temporal associative recall---the \emph{holophone}~\cite{LONGUET_HIGGINS_1968_HOLOPHONE}---showing that linear correlation filters in the frequency domain could store and replay overlapping temporal sequences without dedicated memory slots.

In parallel, connectionist researchers developed matrix-based distributed associative memories, culminating in Hopfield's seminal energy-based attractor network~\cite{HOPFIELD_1982_ASSOCIATIVE}.
A classical Hopfield network stores $N$ bipolar patterns $\{\vect{\xi}_n\}_{n=1}^{N} \subset \{-1, +1\}^D$ inside a symmetric Hebbian outer-product weight matrix $\mat{W} = \frac{1}{D}\sum_{n=1}^{N} \vect{\xi}_n \vect{\xi}_n\trans$ and retrieves corrupted inputs via iterative sign updates minimizing $E(\vect{x}) = -\frac{1}{2}\vect{x}\trans \mat{W} \vect{x}$.
Although Hopfield networks are real-valued and auto-associative rather than complex phase-encoded, their outer-product superposition exhibits the exact crosstalk-limited linear capacity scaling ($N_{\max} \approx 0.138 D$ for uncorrelated patterns) that characterizes all bilinear holographic memories.

\subsubsection{Sutherland's Holographic Neural Network (HNN)}

In the early 1990s, John Sutherland synthesized optical phase interference with neural associative learning in the \emph{Holographic Neural Network (HNN)}~\cite{SUTHERLAND_1990_HNN}.
Rather than mapping scalar features through real-valued hidden layers trained by gradient descent, Sutherland's HNN embeds continuous real-valued stimulus vectors $\vect{s} \in \R^{D_s}$ directly onto the complex unit circle via a smooth sigmoidal phase transformation,
\begin{equation}\label{eq:sutherland_encoding}
    k_j = \lambda_j \exp\!\left(i \theta_j\right), \qquad
    \theta_j = 2\pi \sigma\!\left(\frac{s_j - \mu_j}{\tau_j}\right), \qquad
    \lambda_j \in [0, 1],
\end{equation}
where $\sigma(\cdot)$ denotes the logistic sigmoid (centering the empirical distribution of $s_j$ near zero phase to maximize angular spread across $[-\pi, \pi)$) and $\lambda_j$ encodes stimulus confidence.
Target responses $\vect{v} \in \C^{D_v}$ are phase-encoded analogously.

Sutherland's HNN enfolds a dataset of $N$ stimulus--response associations $\{(\vect{k}_n, \vect{v}_n)\}_{n=1}^N$ into a single complex correlation matrix $\mat{X} \in \C^{D_v \times D_k}$ via non-iterative (one-shot) Hebbian accumulation, $\mat{X} = \sum_{n=1}^{N} \vect{v}_n \vect{k}_n\herm$, as introduced in Remark~\ref{rem:sutherland_matrix}.
When real-world stimuli exhibit non-zero correlations ($\inner{\vect{k}_n}{\vect{k}_q} \neq 0$), one-shot enfolding accumulates systematic bias; Sutherland therefore introduced an iterative error-correcting \emph{enhanced learning} rule (a complex-valued Widrow--Hoff / Delta update),
\begin{equation}\label{eq:sutherland_delta}
    \mat{X}^{(t+1)} = \mat{X}^{(t)} + \eta_t\left(\vect{v}_n - \frac{1}{c_n}\mat{X}^{(t)}\vect{k}_n\right)\vect{k}_n\herm,
    \qquad c_n = \sum_{j=0}^{D_k - 1} \lambda_{n,j},
\end{equation}
which orthogonalizes interfering patterns across passes while remaining strictly local and linear in $\mat{X}$ (reducing at $\eta_0 = 1$, $\mat{X}^{(0)} = \mat{0}$ to one-shot accumulation $\mat{X} = \sum_n \vect{v}_n \vect{k}_n\herm$ with normalized recall $\hat{\vect{v}}_n = c_n^{-1}\mat{X}\vect{k}_n$ matching~\eqref{eq:sutherland_recall} on $\mathbb{T}^{D_k}$, or equivalently placing $c_n^{-1}$ outside the parentheses when $\mat{X}$ absorbs normalization for unnormalized readout $\mat{X}\vect{k}_n$).
To capture nonlinear decision boundaries without deep backpropagation, Sutherland further augmented the input vector with \emph{higher-order phase products}---explicitly constructing composite phasors $\prod_{r \in \mathcal{I}} \lambda_r e^{i\theta_r}$ for index subsets $\mathcal{I}$ of order $p \ge 2$---thereby lifting the stimulus into a higher-dimensional feature space $\C^{D'}$ ($D' \gg D_s$) where both nonlinear expressivity and superposition capacity scale with $D'$.

Despite demonstrating rapid one-shot learning, real-time control, and high fault tolerance in early industrial benchmarks, first-wave HNNs receded from mainstream machine learning during the late 1990s and 2000s.
Their reliance on hand-engineered polynomial phase expansions scaled combinatorially with raw sensory dimensionality, and unlike deep multilayer networks trained with stochastic gradient descent, shallow correlation matrices lacked a mechanism for learning hierarchical, invariant feature representations from raw pixels or natural language.

% ------------------------------------------------------------------------------
\subsection{Second Wave: Vector Symbolic Architectures (HRR and HDC)}\label{sec:history:vsa}
% ------------------------------------------------------------------------------

While Sutherland targeted continuous-valued associative mapping, a parallel community in cognitive science and symbolic AI confronted the \emph{binding problem}: how distributed neural vectors can represent structured, compositional objects---such as nested trees, predicate--argument tuples, and variable bindings---without suffering the superposition catastrophe (conflating ``red square + blue circle'' with ``red circle + blue square'').

\subsubsection{From Tensor Products to Plate's Holographic Reduced Representations}

Paul Smolensky first showed that variable binding in connectionist spaces can be realized cleanly via \emph{Tensor Product Representations (TPRs)}~\cite{SMOLENSKY_1990_TPR}, representing a set of role--filler bindings as $\mat{T} = \sum_{n=1}^{N} \vect{r}_n \otimes \vect{f}_n \in \R^{D \times D}$.
However, recursively binding a composite structure to a higher-level role increases tensor rank at every level of nesting, causing storage complexity to explode exponentially as $\bigo{D^r}$ for depth-$r$ trees.

Tony Plate resolved this dimensionality explosion in 1995 with \emph{Holographic Reduced Representations (HRRs)}~\cite{PLATE_1995_HRR}.
Plate observed that circular convolution $\vect{r} \cconv \vect{f} \in \R^D$ is precisely the \emph{compressed anti-diagonal projection} of the unconjugated outer-product tensor $\vect{r} \otimes \vect{f}$: summing entries $(a,b)$ along wrapped anti-diagonals $(a + b) \bmod D = j$ collapses a rank-2 tensor back into $\R^D$ (whereas summing along wrapped main diagonals $(a - b) \bmod D = j$, as on the Hermitian correlation matrix~\eqref{eq:hnn_corr_matrix}, yields circular correlation $\vect{k} \ccorr \vect{v} = \inv{\vect{k}} \cconv \vect{v}$).
Because the output resides in the exact same ambient space as its operands, HRRs permit arbitrary recursive nesting at constant width $D$: for example, a structured proposition $\texttt{eat}(\texttt{agent}{=}\texttt{dog}, \texttt{object}{=}\texttt{bone})$ can be encoded recursively as
\begin{equation}\label{eq:hrr_recursive}
    \vect{p} = \vect{r}_{\mathrm{pred}} \cconv \vect{f}_{\mathrm{eat}}
             + \vect{r}_{\mathrm{agent}} \cconv \vect{f}_{\mathrm{dog}}
             + \vect{r}_{\mathrm{obj}} \cconv \vect{f}_{\mathrm{bone}} \in \R^D,
\end{equation}
and subsequently bound as a filler inside an outer sentence vector without expanding dimensionality.
Plate further formalized the clean-up memory required to strip crosstalk noise after circular-correlation unbinding ($\ccorr$), proved the foundational variance bounds analyzed in Section~\ref{sec:math:capacity}, and introduced complex-valued \emph{Fourier HRR (FHRR)} on $\mathbb{T}^D$ to eliminate spectral self-distortion and accelerate binding to $\bigo{D}$ elementwise phase addition.
More recently, Generalized HRRs (GHRRs)~\cite{YEUNG_2024_GHRR} lifted scalar FHRR phasors in $\mathbb{T}^D$ to $m \times m$ unitary matrix elements in $U(m)^D$, enabling natively non-commutative binding and flexible similarity kernels while recovering standard FHRR at $m=1$.
Bridging these second-wave algebras directly into deep sequence modeling, Alam et al.~\cite{HOLOGRAPHIC_ATTENTION_OR_KV_CACHE} recast Transformer self-attention via compressive HRR binding and unbinding to reduce attention complexity to linear-in-$N$ superposition traces.

\subsubsection{Kanerva's Hyperdimensional Computing and VSA Unification}

Building on sparse distributed memory and random-projection geometry, Pentti Kanerva articulated \emph{Hyperdimensional Computing (HDC)}~\cite{KANERVA_2009_HDC}, emphasizing ultra-high-dimensional ($D \sim 10^4$) random binary or bipolar \emph{hypervectors} ($\{0, 1\}^D$ or $\{-1, +1\}^D$).
In addition to elementwise binding ($\hadamard$ or XOR) and majority-vote/additive bundling, Kanerva highlighted a third primitive---invertible coordinate \emph{permutation} $\rho: \R^D \to \R^D$ (typically a cyclic shift)---to break the commutativity of binding and encode ordered sequences, $n$-grams, and temporal trajectories:
\begin{equation}\label{eq:hdc_sequence}
    \vect{s}_{1:L} = \sum_{t=1}^{L} \rho^{L - t}\!\left(\vect{x}_t\right)
    \qquad\text{or}\qquad
    \vect{g}_{1:L} = \bigodot_{t=1}^{L} \rho^{L - t}\!\left(\vect{x}_t\right).
\end{equation}
Kanerva demonstrated that holistic algebraic analogies (such as solving ``What is the Dollar of Mexico?'' by unbinding $\vect{k}_{\mathrm{USA}}$ and rebinding $\vect{k}_{\mathrm{Mexico}}$ across bundled country records) execute in a single constant-time parallel vector operation without explicit symbolic search trees.

Subsequent surveys by Schlegel et al.~\cite{SCHLEGEL_VSA_COMPARISON} and Kleyko et al.~\cite{KLEYKO_VSA_SURVEY} unified HRR, FHRR, HDC, Multiply--Add--Permute (MAP), Binary Spatter Codes (BSC), and Sparse Binary Distributed Representations under the umbrella term \emph{Vector Symbolic Architectures (VSAs)}.
Across this taxonomy, every VSA is uniquely characterized by four design choices: (1) ambient vector space and atomic distribution ($\{0,1\}^D$, $\{-1,+1\}^D$, $\R^D$, or $\mathbb{T}^D$); (2) binding algebra and whether unbinding is exact (e.g., XOR, MAP, FHRR) or approximate (real Gaussian HRR); (3) bundling linearity vs.\ clipping/phase-projection; and (4) clean-up mechanism (linear codebook matching, Hopfield attractors, or multi-factor resonator networks).

% ------------------------------------------------------------------------------
\subsection{Third Wave: Spectral Mixing, Complex Deep Learning, and State-Space Models}\label{sec:history:modern}
% ------------------------------------------------------------------------------

Over the past decade, deep learning architectures designed for scalability and long-sequence stability have independently rediscovered---and scaled via end-to-end backpropagation---each core mathematical ingredient of holographic computation.

\subsubsection{Complex-Valued Neural Networks (CVNNs) and Unitary Recurrences}

Modern Complex-Valued Neural Networks~\cite{BASSEY_CVNN_SURVEY,TRABELSI_2018_DCN} extend deep convolutional, residual, and attention layers natively into $\C^D$.
Because Liouville's theorem prohibits non-constant bounded holomorphic activations on $\C$, CVNNs rely on \emph{Wirtinger ($\C\R$) calculus} to optimize real-valued objective functions $\mathcal{L}(\vect{z}, \conj{\vect{z}}) \in \R$ with respect to complex weights $\vect{z} = \vect{x} + i\vect{y}$ via the conjugate gradient
\begin{equation}\label{eq:wirtinger}
    \nabla_{\conj{\vect{z}}} \mathcal{L}
    \coloneqq \frac{\partial \mathcal{L}}{\partial \conj{\vect{z}}}
    = \frac{1}{2}\left(\frac{\partial \mathcal{L}}{\partial \vect{x}} + i\frac{\partial \mathcal{L}}{\partial \vect{y}}\right),
\end{equation}
paired with split or phase-preserving nonlinearities such as $\C$ReLU (applying ReLU separately to real and imaginary parts) and $\mathrm{modReLU}(z) = \mathrm{ReLU}(\abs{z} + b)\frac{z}{\abs{z}}$ (which preserves phase angle $\theta$ identically while gating magnitude $\lambda$).
Closely aligned with unitary holographic binding, \emph{Unitary Evolution Recurrent Neural Networks (uRNNs)}~\cite{ARJOVSKY_2016_URNN} parameterize recurrent transitions as products of complex diagonal phase matrices, FFT/IFFT blocks, and Householder reflections so that $\mat{U}\herm \mat{U} = \mat{I}_D$.
Because unitary operators preserve Euclidean norm exactly ($\norm{\mat{U}\vect{h}}_2 = \norm{\vect{h}}_2$), uRNNs eliminate exploding and vanishing gradients over long horizons purely through complex phase interference---bridging Sutherland's phase-space formulation with deep recurrent optimization.

\subsubsection{Fourier Token Mixing and Spectral Convolutions}

A second point of convergence emerged in efforts to replace quadratic $\bigo{N^2}$ Transformer self-attention~\cite{VASWANI_TRANSFORMER} with sub-quadratic spectral mixing.
FNet~\cite{LEE_THORP_FNET} replaces the multi-head self-attention sublayer entirely with an unparameterized two-dimensional Discrete Fourier Transform across the sequence and hidden dimensions,
\begin{equation}\label{eq:fnet}
    \mat{Y} = \Real\!\left(\DFT_{\mathrm{seq}}\!\left(\DFT_{\mathrm{hidden}}(\mat{H})\right)\right),
\end{equation}
achieving $\bigo{N \log N}$ token mixing while retaining competitive downstream accuracy on natural-language benchmarks at a fraction of the computational cost.
Subsequent global-filter and FFT long-convolution models generalize~\eqref{eq:fnet} by introducing learnable frequency-domain filters $\vect{W}_\omega \in \C^D$, computing $\IDFT(\vect{W} \hadamard \DFT(\vect{x}))$ across sequence positions.
By the Circular Convolution Theorem (Proposition~\ref{prop:conv_thm}), every such spectral mixing layer is algebraically equivalent to circular-convolution holographic binding between input representations and a learned kernel, executed in $\bigo{N \log N}$ time via FFT.

\subsubsection{Linear Attention, State-Space Models, and Modern Hopfield Networks}

Most remarkably, modern sub-quadratic sequence models have converged directly onto the exact mathematical form of holographic associative superposition.
Katharopoulos et al.~\cite{KATHAROPOULOS_2020_LINEAR_TRANSFORMER} showed that replacing the exponential softmax kernel in Transformer attention with a separable positive feature map $\phi(\cdot)$ rewrites causal attention as a recurrent fixed-size associative matrix,
\begin{equation}\label{eq:linear_attn_state}
    \mat{S}_t = \mat{S}_{t-1} + \vect{v}_t \phi(\vect{k}_t)\trans,
    \qquad
    \hat{\vect{v}}_t = \frac{\mat{S}_t \phi(\vect{q}_t)}{\vect{z}_t\trans \phi(\vect{q}_t)},
\end{equation}
where $\vect{z}_t = \sum_{\tau=1}^{t} \phi(\vect{k}_\tau)$ normalizes the readout.
Comparing~\eqref{eq:linear_attn_state} with Sutherland's complex correlation matrix~\eqref{eq:sutherland_recall} reveals an exact structural isomorphism: linear attention is a real-valued, end-to-end learned holographic outer-product memory whose fixed state size $\bigo{D^2}$ suffers precisely from the crosstalk saturation analyzed in Section~\ref{sec:math:capacity}.

Structured and Selective State-Space Models (S4 and Mamba)~\cite{GU_SSM_S4_MAMBA} resolve the static saturation of unbounded additive traces by introducing continuous-time (frequently complex-diagonalized) dynamics and input-dependent forgetting gates:
\begin{equation}\label{eq:ssm_recurrence}
    \mat{S}_t = \mat{A}_t \mat{S}_{t-1} + \vect{B}_t \vect{x}_t,
    \qquad
    \vect{y}_t = \mat{C}_t \mat{S}_t,
\end{equation}
where $\mat{A}_t$ selectively decays stale superposed bindings to preserve $\SNR$ for salient recent or key-matched information---the continuous, learned counterpart to weighted holographic bundling~\eqref{eq:bundling} and Sutherland's delta-rule interference suppression~\eqref{eq:sutherland_delta}.
Extending this matrix-state lineage, DeltaNet~\cite{YANG_2024_DELTANET} and Gated DeltaNet~\cite{YANG_2025_GATED_DELTANET} scale Sutherland's error-correcting delta rule~\eqref{eq:sutherland_delta} into chunkwise hardware-efficient linear Transformers, combining rank-one key--value erasure with selective decay to actively suppress cross-key interference.
In parallel, frontier hybrid models---exemplified by Griffin~\cite{DE_2024_GRIFFIN} (interleaving gated linear recurrences with local sliding-window attention) and Titans~\cite{BEHROUZ_2025_TITANS} (pairing local exact attention with a deep neural long-term memory updated at test time via surprise-driven gradient descent)---operationalize the exact division of labor between fixed-state compression and exact local recall.
In stark contrast to this fixed-size parametric superposition lineage, Ramsauer et al.~\cite{RAMSAUER_MODERN_HOPFIELD} proved that lifting Hopfield's classical quadratic energy to a continuous log-sum-exp energy yields an associative memory with \emph{exponential} pattern storage capacity whose single-step update is mathematically identical to Transformer scaled dot-product softmax attention~\cite{VASWANI_TRANSFORMER}---at the unavoidable cost of retaining every stored pattern explicitly in $\bigo{ND}$ memory.

% ------------------------------------------------------------------------------
\subsection{Unifying Architectural Perspective}\label{sec:history:synthesis}
% ------------------------------------------------------------------------------

Viewing all three waves through the mathematical lens of Section~\ref{sec:math} exposes a clean architectural bifurcation across modern AI:
\begin{itemize}
    \item \emph{Parametric holographic superposition} (Sutherland HNNs, HRR/FHRR, HDC, linear attention, and SSMs) compresses unbounded histories into a constant-size latent state ($\bigo{D}$ for vector binding or $\bigo{D^2}$ for outer-product matrices), delivering $\bigo{1}$ per-step inference memory and native phase/hardware affinity at the cost of additive crosstalk governed by $\SNR \propto D/N$.
    \item \emph{Non-parametric explicit attention} (Modern Hopfield Networks and softmax Transformers) avoids crosstalk by retaining all $N$ keys and values explicitly ($\bigo{ND}$ memory and $\bigo{N^2}$ pairwise interaction), enabling exact verbatim recall at severe long-context memory and energy overhead.
\end{itemize}
This dichotomy forms the analytical foundation for the systematic multi-axis comparison in Section~\ref{sec:comparison} and motivates the hybrid holographic--attention architectures examined in Section~\ref{sec:potentials}.
